\section{Exact pair-key geometry}

Throughout the following finite-field arithmetic sections, p is an odd prime,
$H\leq\mathbb F_p^\times$ has odd prime order K, and
$e_p(t)=e^{2\pi i t/p}$. Fixed-degree row presentations impose $K\geq7$
where stated. These hypotheses are distinct from the composite Pell moduli
in the later quantum sections.

For two admissible position pairs $(x_i,y_i)\in(\mathbb F_p\setminus\{0,1\})^2$,
write their keys as $(u_i,v_i,d_i)$. Let $B(t)$ count the ordered quadruples
satisfying $u_1u_2\in H$, $v_1v_2\in H$, and $d_1+d_2=t$; let $O(t)$ count
those for which both position pairs are off diagonal. Write $V_4$ for the
set of all such quadruples without fixing t. These definitions include
multiplicities in the collapsed diagonal fibre.

For $x,y\in\Fp\setminus\{0,1\}$ define
\begin{equation}\label{eq:keydef}
 u=\frac{x}{y},\qquad
 v=\frac{x-1}{y-1},\qquad
 d=x-y.
\end{equation}

\begin{proposition}[Pair-key bijection]\label{prop:key}
If $x\ne y$, then $u,v\in\Fp^\times\setminus\{1\}$, $u\ne v$, and
\begin{equation}\label{eq:inverse}
 y=\frac{1-v}{u-v},\qquad
 x=\frac{u(1-v)}{u-v},\qquad
 d=\frac{(u-1)(v-1)}{v-u}.
\end{equation}
Conversely every such $(u,v)$ gives a unique admissible off-diagonal pair.
The diagonal collapses to $(u,v,d)=(1,1,0)$ with multiplicity $p-2$.
\end{proposition}

\begin{proof}
The equations $x=uy$ and $x-1=v(y-1)$ give $(u-v)y=1-v$, hence
\eqref{eq:inverse}.  The diagonal statement is immediate.
\end{proof}

Under the subgroup constraints $u,v\in H$, there are exactly
$(K-1)(K-2)$ admissible off-diagonal ordered pairs.  If $O(t)$ denotes the
off/off contribution and
\[
C_H(t)=\#\left\{(u,v)\in(H\setminus\{1\})^2:
 u\ne v,\ \frac{(u-1)(v-1)}{v-u}=t\right\},
\]
then exact bookkeeping gives
\begin{equation}\label{eq:Bdecomp}
B(t)=O(t)+2(p-2)C_H(t)+(p-2)^2\mathbf1_{t=0}.
\end{equation}
The factor two is essential: diagonal/off and off/diagonal are distinct ordered
classes.  The total count is
\begin{equation}\label{eq:V4}
|V_4|=(K^2+1)p^2-(6K^2+3K)p+13K^2-3K.
\end{equation}
If $\Sigma_4(q)=\sum_tB(t)e_p(qt)$, $V=\sum_tB(t)$ and $Z=B(0)$, finite
Fourier orthogonality gives
\begin{equation}\label{eq:meanvar}
\mathbb E_{q\ne0}\Sigma_4(q)=\frac{pZ-V}{p-1},
\qquad
\Var_{q\ne0}\Sigma_4(q)=p\Var_{t\ne0}B(t).
\end{equation}
No probabilistic independence assumption enters this reduction.

\section{Cayley energy and the best uniform bound}

Put
\[
c(h)=\frac{h+1}{h-1},\qquad A=c(H\setminus\{1\}).
\]
For $u\ne v$,
\begin{equation}\label{eq:cayley}
\frac{(u-1)(v-1)}{v-u}=\frac{2}{c(u)-c(v)},
\qquad c(h^{-1})=-c(h).
\end{equation}
Thus $A=-A$ and the nonzero multiplicities of $C_H$ are additive-difference
multiplicities of $A$ up to the permutation $x\mapsto2/x$.  Hence
\begin{equation}\label{eq:energyCH}
\sum_tC_H(t)^2=E_+(A)-|A|^2,
\qquad
E_+(A)=\sum_xr_{A-A}(x)^2.
\end{equation}
Identity and central inversion force the nonzero baseline
\begin{equation}\label{eq:forced}
(K-1)(2K-5).
\end{equation}
For the unnormalized Fourier transform,
\begin{equation}\label{eq:centered}
E_+(A)=\frac{(K-1)^4}{p}
 +\frac1p\sum_{q\ne0}|\widehat{1_A}(q)|^4.
\end{equation}
The sharp pointwise target is therefore
\begin{equation}\label{eq:sharp-target}
E_+(A)=\frac{(K-1)^4}{p}+O(K^{2+o(1)}).
\end{equation}

For $r_A(t)=\#\{(a,b)\in A^2:a+b=t\}$, clearing Cayley denominators gives
\[
(2-t)xy+t(x+y)-(2+t)=0,
\]
an absolutely irreducible curve in odd characteristic: when $t\ne2$ the
bilinear coefficient determinant is $-4\ne0$, and when $t=2$ it is a
nonconstant linear polynomial. Its bidegree is $(1,1)$ in either case.
For $t\ne0,-2$ the constant term is nonzero and $\deg P(x,0)=1$, exactly
the hypotheses needed for the subgroup-curve input.

\begin{theorem}[Direct pointwise Cayley bound]\label{thm:direct}
If
\[
100<K<\frac13p^{3/4},
\]
then
\begin{equation}\label{eq:direct}
E_+(A)<34K^{8/3}.
\end{equation}
\end{theorem}

\begin{proof}
For $t\ne0,-2$, the source form of Vyugin--Makarychev's subgroup-curve theorem
\cite{VM2015,MV2019} gives $r_A(t)\le32K^{2/3}$ after its degree hypotheses are
checked on the displayed bilinear polynomial.  The exceptional fibres satisfy
$r_A(0)=K-1$ and $r_A(-2)\le K-1$.  Since
$\sum_tr_A(t)=(K-1)^2$,
\[
E_+(A)
\le2(K-1)^2+32K^{2/3}\sum_{t\ne0,-2}r_A(t)
<34K^{8/3}.
\]
\end{proof}

\section{Prime-order rigidity, square resultants and discriminants}

Let $\zeta=e^{2\pi i/K}$.  Then
\[
c(\zeta^a)=-i\cot\frac{\pi a}{K}.
\]
Chowla's prime-modulus cotangent independence \cite{Chowla1970,Wang1982}
identifies the characteristic-zero Cayley set with a rational cross-polytope.

\begin{theorem}[Prime-order additive rigidity]\label{thm:add-rigidity}
Every characteristic-zero equal-difference collision among
$\{c(\zeta^a):a\in\mathbf F_K^*\}$ is identity or central inversion-swap.
\end{theorem}

For every admissible cleared collision polynomial $P(X)$, direct algebra gives
$P(X)=(X-1)^3Q(X)$ and a reciprocal relation for $Q$.  Passing to the maximal
real subfield yields:

\begin{theorem}[Square-resultant theorem]\label{thm:square-res}
\begin{equation}\label{eq:square-res}
\Res(\Phi_K,P)=K^3M^2,\qquad M\in\Z.
\end{equation}
For a non-forced characteristic-zero collision, $M\ne0$.
\end{theorem}

\begin{proof}
Here the cleared equal-difference polynomial can be defined explicitly.
For $1\leq a_i<K$ and $(\epsilon_1,\epsilon_2,\epsilon_3,\epsilon_4)
=(1,-1,-1,1)$, put
\[
 P(X)=\sum_{i=1}^4\epsilon_i(X^{a_i}+1)
                         \prod_{j\ne i}(X^{a_j}-1),\qquad s=\sum_i a_i.
\]
Every term contains three factors vanishing at one, so
$P=(X-1)^3Q$ with $Q\in\mathbb Z[X]$. Direct inversion gives
$P(X^{-1})=-X^{-s}P(X)$ and hence $Q(X^{-1})=X^{3-s}Q(X)$.
Choose an integer k with $2k\equiv s-3\pmod K$. Then
$\xi=\zeta^{-k}Q(\zeta)$ is real and is an algebraic integer.
The full cyclotomic norm of Q is the square of the real-subfield norm
of $\xi$, because $\operatorname{Norm}(\zeta)=1$.
Also $\operatorname{Norm}(\zeta-1)=K$. Therefore
\[
 \operatorname{Res}(\Phi_K,P)
       =K^3\operatorname{Norm}_{\mathbb Q(\zeta+\zeta^{-1})/\mathbb Q}(\xi)^2.
\]
The inner norm is an integer. In the non-forced case the additive rigidity
